\clearpage
\setcounter{aufgabe}{7}
\einheitskopf{Einheit 2 von 5 · Prozentsatz berechnen}\label{e2}

\begin{aufgabe}{Prozentsatz – Was ist das Ganze? Unterstreiche in jedem Text das Ganze und schreibe es in das Feld. Rechne noch nicht.}
\begin{teile}
\teil In einer Klasse mit 25 Kindern tragen 10 eine Brille. \hfill Das Ganze: \leerfeld
\teil Von 60 Zuschauern gehen 15 in der Pause nach Hause. \hfill Das Ganze: \leerfeld
\teil Ein Fahrrad kostete $400\,€$, jetzt kostet es $340\,€$. \hfill Das Ganze: \leerfeld
\teil In einem Korb liegen 8 grüne und 12 rote Äpfel. \hfill Das Ganze: \leerfeld
\teil Von den 30 Losen sind 6 Gewinne, der Rest sind Nieten. \hfill Das Ganze: \leerfeld
\end{teile}
\end{aufgabe}

\begin{aufgabe}{Prozentsatz – am Streifen. Der ganze Streifen ist das Ganze. Wie viel Prozent des Ganzen ist der graue Teil? Beispiel: Sind 2 der 10 Kästchen grau, so sind das 20\,\%.}
\streifenwertreihe{a/30/6 kg/20 kg, b/70/35 €/50 €, c/40/32 Schüler/80 Schüler, d/45/90 m/200 m}
\begin{teilezwei}
\tz \leerfeld[\%] & \tz \leerfeld[\%] \\
\tz \leerfeld[\%] & \tz \leerfeld[\%] \\
\end{teilezwei}
\end{aufgabe}

\begin{aufgabe}{Prozentsatz – rechnen. Teile den Teil durch das Ganze und schreibe das Ergebnis als Prozentangabe.}
\beispiel{19 \text{ von } 100: \quad 19 : 100 = 0{,}19 = 19\,\%}
\begin{teilezwei}
\tz 19 von 100: \leerfeld[\%] & \tz 7 von 100: \leerfeld[\%] \\
\tz 64 von 100: \leerfeld[\%] & \tz 90 von 100: \leerfeld[\%] \\
\tz 21 von 25: \leerfeld[\%] & \tz 21 von 60: \leerfeld[\%] \\
\end{teilezwei}
Runde auf eine Stelle nach dem Komma.
\begin{teile}
\teil 5 von 12: \leerfeld[\%]
\end{teile}
\end{aufgabe}

\begin{aufgabe}{Prozentsatz – Umkehrung. Denke dir ein Zahlenpaar aus, das passt.}
\begin{teile}
\teil Der Teil soll $50\,\%$ des Ganzen sein. \hfill Teil: \leerfeld \ Ganzes: \leerfeld
\teil Der Teil soll $25\,\%$ des Ganzen sein. \hfill Teil: \leerfeld \ Ganzes: \leerfeld
\teil Gib zwei verschiedene Zahlenpaare an, bei denen der Teil $40\,\%$ des Ganzen ist. \\ Teil: \leerfeld \ Ganzes: \leerfeld \\ Teil: \leerfeld \ Ganzes: \leerfeld
\end{teile}
\end{aufgabe}

\begin{aufgabe}{Prozentsatz – aus dem Sachtext. Suche zuerst das Ganze, dann rechne.}
\begin{teile}
\teil In einem Kinosaal sind 40 Sitzplätze, 26 davon sind besetzt. Wie viel Prozent der Plätze sind besetzt? \leerfeld[\%]
\teil Ein Verein hat 18 weibliche und 32 männliche Mitglieder. Wie viel Prozent der Mitglieder sind weiblich? \leerfeld[\%]
\teil Eine Jacke kostete $80\,€$, jetzt kostet sie $68\,€$. Um wie viel Prozent ist sie billiger geworden? \leerfeld[\%]
\teil In einer Schale liegen 7 grüne und 17 rote Äpfel. Wie viel Prozent der Äpfel sind grün? Runde auf eine Stelle nach dem Komma. \hfill (P10 2018) \\ \leerfeld[\%]
\end{teile}
\end{aufgabe}

\begin{aufgabe}{Prozentsatz – Fehler finden. In ein Regal passen 40 Bücher, 24 stehen darin. Nina will wissen, wie viel Prozent des Regals gefüllt sind, und rechnet:}
\rechnung{40 : 24 &= 1{,}666\ldots \approx 166{,}7\,\%}
\begin{teile}
\teil Benenne den Fehler und korrigiere die Rechnung. \\ Fehler: \dsleer\dsleer\dsleer \quad richtig: \leerfeld[\%]
\teil In einem Bus gibt es 45 Sitzplätze, 36 sind besetzt. Wie viel Prozent der Plätze sind besetzt? \leerfeld[\%]
\end{teile}
\end{aufgabe}

\begin{aufgabe}{Prozentsatz – Begründen. Begründe in einem Satz.}
\begin{teile}
\teil Warum teilt man den Teil durch das Ganze und nicht das Ganze durch den Teil? \\ \dsleer\dsleer\dsleer\dsleer
\teil Tom sagt: „Wenn beim Prozentsatz mehr als $100\,\%$ herauskommen, habe ich mich bei einem Teil vom Ganzen verrechnet.“ Hat er recht? \janein \\ Begründung: \dsleer\dsleer\dsleer\dsleer
\end{teile}
\end{aufgabe}

\begin{aufgabe}{Prozentsatz – Anwendung. Beim Basketballtraining wirft Ali 20-mal auf den Korb und trifft 12-mal. Ben wirft 35-mal und trifft 21-mal.}
\begin{teile}
\teil Wie viel Prozent seiner Würfe trifft Ali? \leerfeld[\%]
\teil Wie viel Prozent seiner Würfe trifft Ben? \leerfeld[\%]
\teil Wer hat die bessere Trefferquote? Begründe. \\ \dsleer\dsleer\dsleer\dsleer
\end{teile}
\end{aufgabe}
